\section{Compare five clipped scenes}\label{blit-x-trace}
\begin{table}[htbp]
  \centering
  \caption{Expected counts for the supplied five-by-four sprite and eight-by-six framebuffer. The key is hexadecimal 00ff00ff.}
  \tagpdfsetup{table/header-rows={1}}
  \begin{tabular}{rrrrrr}
    \toprule
    $x$  & $y$  & Copy width & Copy height & Opaque \texttt{v0} & Visible \texttt{v1} \\
    \midrule
    $-2$ & 1    & 3          & 4           & 8                  & 12                  \\
    0    & 0    & 5          & 4           & 13                 & 20                  \\
    6    & $-2$ & 2          & 2           & 2                  & 4                   \\
    $-1$ & 4    & 4          & 2           & 6                  & 8                   \\
    8    & 0    & 0          & 4           & 0                  & 0                   \\
    \bottomrule
  \end{tabular}
\end{table}
For the top-right clip, source rows zero and one lie above the framebuffer,
and source columns two through four lie to its right. Only source rows two
and three, columns zero and one remain. They contain a key, white, red, and
a key. The two opaque values explain the result without relying on the display.

For exact contact at $x=8$, horizontal width is zero. No row loop runs, even
though the vertical intervals overlap. Both returned counts are zero, and no
pixel load or store is needed. Check this case before measuring an optimization
that attempts to skip completely invisible sprites.
